%% Abstract - owner: Taj. Written last, once the results existed.
%% 150-200 words. Every number here is measured and reproducible:
%%   make -C sim test            8/8 testbenches
%%   make -C sim system          golden-model comparison
%%   make -C sim scalar_top      64 cycles
%%   python synth/run_synth.py   area figures
%%   python fuzz/run_fuzz.py     random-program equivalence

Data-level parallelism is the cheapest form of parallelism to exploit in
hardware, but the gain it offers is easy to assert and hard to substantiate.
This paper presents a parameterised four-lane SIMD vector unit implemented in
Verilog, together with the scalar baseline and verification apparatus needed to
measure what the vector hardware actually buys. The unit provides eight 128-bit
vector registers and a frozen 16-bit instruction set covering vector addition,
subtraction, bitwise operations and signed comparison, with the lane count held
as a synthesis-time parameter. Correctness is established against an independent
Python model of the instruction set rather than against hand-written
expectations: a self-checking system testbench compares every result word, and a
random-program fuzzer runs generated programs on both the hardware and the model
and compares memory and cycle counts. That fuzzer exposed a reset-time
write-enable defect that eight directed testbenches had missed. On a 16-element
array addition the vector unit completes in 16 cycles against 64 for a scalar
baseline built from the same ALU, a speed-up of $4.00\times$ for a
$4.03\times$ increase in LUT count.
